\subsection{Interfaz, trama de metadatos y votación}

El programa agrega al audio los metadatos de entrada, de haber alguno, y escribe el resultado en un archivo WAV de salida; de otra forma, se limita a imprimir los metadatos que encuentre.

Como ilustra la Tabla~\ref{tab:registro}, cada metadato entrante se convierte en un registro con su texto en el formato de transformación de Unicode de 8 bits (UTF-8)~\cite{rfc3629}, de 1 a 255 bytes, y una verificación de redundancia cíclica (CRC) de 16 bits que permite detectar si algún bit del registro llegó mal. La CRC usa el polinomio $x^{16}+x^{12}+x^{5}+1$, con valor inicial \texttt{0xFFFF} y sin reflexión ni operación final~\cite{reveng_crc}. Cada registro tiene un campo para su categoría con la siguiente taxonomía en orden de 1 a 7: título, artista, autor, álbum, año, artista invitado y productor. Los registros se concatenan en orden de categoría, con el bit más significativo de cada byte primero, y forman el mensaje. Cada ventana de 1650 muestras lleva un bit de información, y el mensaje que genera se repite hasta ocupar todas las ventanas. Si el audio no alcanza para una copia completa, se codifica lo que cabe con un aviso. Las ventanas con todas sus muestras en cero no llevan bit, porque siguen en cero después de agregarles el eco; el codificador las salta, como proponen Gruhl \textit{et al.}~\cite{gruhl_echo_hiding}.

\begin{table}[!t]
\caption{Campos de un registro de la trama de metadatos.}
\label{tab:registro}
\centering
\begin{tabular}{lcl}
\toprule
Campo & Bits & Contenido \\
\midrule
Preámbulo & 8 & \texttt{0xB2} \\
Categoría & 8 & 1 a 7 \\
Longitud & 8 & $\ell$, de 1 a 255 bytes \\
Datos & $8\ell$ & texto en UTF-8 \\
CRC & 16 & sobre categoría, longitud y datos \\
\bottomrule
\end{tabular}
\end{table}

El decodificador tolera un bit errado en el preámbulo cuando está justo donde se espera que empiece un registro, al inicio o justo después de uno leído bien; lo mismo no es cierto para otras posiciones, ya que se podría confundir con un byte de datos. Si la CRC coincide, el registro es válido cuando su categoría está entre 1 y 7 y su texto es UTF-8 correcto, y desconocido en otro caso; si no coincide y el registro empezaba donde se esperaba, se considera dañado. Si un registro no es confirmado, la búsqueda sigue desde el bit siguiente en vez de saltar el registro entero: así, si se pierde o se agrega un bit, solo se pierde ese registro y los siguientes se vuelven a encontrar.

Si cada bit se lee mal con probabilidad $p$, de forma independiente, un registro de $m$ bits llega intacto con probabilidad $(1-p)^m$: con $p = \SI{5}{\percent}$, \SI{2.5}{\percent} para un registro de 72 bits y \SI{0.0007}{\percent} para uno de 232. Por eso se agregó un mecanismo de votación para la corrección de errores: las copias repetidas del mensaje se combinan con un código de repetición con decisión por mayoría. Para votar hay que saber dónde empieza cada copia. Como el mensaje se repite, la secuencia de $B$ bits leída se parece a sí misma desplazada un período $P$: se prueba cada desplazamiento, de 48 a $\min(\lfloor B/3\rfloor,14560)$ bits, y se cuenta cuántos bits coinciden, en bloques de 64 bits que aceptan hasta 4 bits de corrimiento. Como $2P$ y $3P$ también coinciden, se elige el menor desplazamiento a menos de 2 puntos porcentuales del mejor, y se ajusta probando sin tolerancia ese valor y los 8 siguientes; si ni el mejor llega al \SI{75}{\percent} de bits iguales, no se vota. Superponiendo las copias, cada bit se decide por mayoría; en un empate, o con menos de tres copias, vale la primera. Después, con la copia votada como referencia, se vuelve a buscar el inicio de cada copia, hasta 4 bits antes o después, y se vota otra vez. Este proceso se repite dos veces. La copia votada se lee igual que la secuencia original, y sus registros válidos se suman a los de la lectura directa.

Por cada categoría con algún registro válido se imprime en la salida estándar su nombre y el texto que más se repite (por ejemplo, «Título: The Entertainer»), y los registros desconocidos o dañados se imprimen en la salida de error con una advertencia, solo si se leyó algún registro válido, para no dar avisos falsos con un audio sin codificar.
